% capitulo2.tex
\chapter{Planteamiento y Objetivos de la Investigación}

% ============================================================
% SECCIÓN 2.1 — PLANTEAMIENTO DEL PROBLEMA
% ============================================================
\section{Planteamiento del Problema}

En el ámbito de la ingeniería estructural y la sismorresistencia, la conexión entre la
superestructura metálica y la cimentación de concreto reforzado representa el eslabón más crítico
de la cadena resistente de un edificio de acero. Las \textit{placas base} y sus sistemas de
anclaje fungen como la interfaz fundamental encargada de transferir las fuerzas axiales $P_u$,
los esfuerzos cortantes basales $V_u$ y los momentos flexionantes $M_u$ generados en las columnas
hacia los pedestales y zapatas de concreto \cite{FisherKloiber_2006}. La integridad de esta
conexión es determinante para la respuesta global de la estructura ante eventos sísmicos severos,
ya que su falla ---particularmente cuando ocurre de manera frágil y repentina--- puede precipitar
mecanismos de colapso que comprometan la seguridad humana y la funcionalidad de la edificación
\cite{AISC_341_2022}.

\begin{figure}[H]
  \centering
  \includegraphics[width=0.55\textwidth]{\detokenize{figuras/001_Vista en corte de una conexión columna-cimentación columna de acero, placa base, pernos de anclaje, grout y pedestal de concreto..png}}
  \caption{Vista en corte transversal de los componentes de la conexión columna-cimentación:
    columna de acero, placa base, pernos de anclaje, mortero de nivelación (grout) y pedestal de
    concreto \cite{GarciaPalencia_2008}.}
  \label{fig:corte_conexion}
\end{figure}

La República Mexicana se asienta sobre la convergencia de cinco placas tectónicas mayores
---Norteamérica, Caribe, Cocos, Rivera y Pacífico---, lo que la convierte en uno de los países con
mayor amenaza sísmica del planeta. La Ciudad de México concentra además una condición agravante:
sus zonas lacustres están constituidas por arcillas altamente compresibles que amplifican
dinámicamente el movimiento del suelo, incrementando las aceleraciones espectrales de forma
sustancial respecto a sitios de roca firme \cite{UNAM_SSN_2017}. Esta particularidad geológica
eleva la demanda sobre todos los componentes estructurales, y de manera especialmente crítica
sobre las conexiones columna-cimentación, que son las primeras en acusar las solicitaciones de
mayor amplitud.

Durante décadas, el diseño de placas base en México se realizó bajo filosofías simplificadas que
modelaban la placa como un elemento rígido con distribución lineal de presiones de contacto sobre
el concreto \cite{GCDMX_NTC_DCEA_2017}. Este enfoque omitía dos fenómenos fundamentales: la
flexibilidad real de la placa ---que puede incrementar la demanda en los pernos de anclaje mediante
el efecto de palanca conocido como \textit{prying action}--- y la naturaleza frágil del mecanismo
de falla por arrancamiento de cono de concreto (\textit{concrete breakout}), que puede precipitar
el colapso de la conexión sin advertencia previa \cite{FisherKloiber_2006}.

El 19 de septiembre de 2017, el sismo de magnitud momento $M_w~7.1$ con epicentro en la región
Puebla-Morelos demostró con crudeza las insuficiencias del paradigma normativo vigente. La
Figura~\ref{fig:mapa_sismo2017} ilustra la proximidad del evento con la capital. Las inspecciones
forenses documentadas por la Sociedad Mexicana de Ingeniería Estructural (\textsc{SMIE})
identificaron, de manera reiterada, fallas en conexiones columna-base de estructuras de acero
diseñadas conforme a las disposiciones entonces vigentes \cite{SMIE_2017}. El modo de falla
predominante fue el arrancamiento frágil del cono de concreto en los pernos de anclaje, que
colapsó antes de que la columna metálica pudiera fluir dúctilmente para disipar la energía
sísmica \cite{UNAM_SSN_2017}.

\begin{figure}[H]
  \centering
  \includegraphics[width=0.62\textwidth]{\detokenize{figuras/002_Mapa de la República Mexicana señalando el epicentro del sismo del 19-S 2017 y la Ciudad de México..png}}
  \caption{Ubicación geográfica del epicentro del sismo intraplaca del 19 de septiembre de 2017
    ($M_w~7.1$), ilustrando su cercanía con la Ciudad de México y las principales zonas urbanas
    afectadas \cite{SSN_Puebla_2017}.}
  \label{fig:mapa_sismo2017}
\end{figure}

La revisión normativa que derivó en las Normas Técnicas Complementarias para el Diseño y
Construcción de Estructuras de Acero 2023 (NTC-DCEA 2023) introdujo cambios sustanciales: la
metodología de diseño de placas base fue promovida del Apéndice~B de la edición 2017 al Capítulo
13 formal de la norma actualizada. Se incorporaron dos nuevos factores de confinamiento por
refuerzo transversal del dado de concreto: el factor $R_{et}$ para resistencia al cono de tensión
y el factor $R_{ev}$ para resistencia al cono por cortante \cite{GCDMX_NTC_DCEA_2023}. Además, la
distancia mínima del perno al borde libre del pedestal fue modificada de $6\,d_o$ a $7\,d_o$
para aceros de alta resistencia A325/A449.

Estos cambios plantean una brecha de conocimiento técnico de relevancia inmediata para la práctica
estructural en México: se desconoce de qué forma y en qué magnitud la transición entre ambas
normas modifica cuantitativamente los parámetros de diseño del sistema
placa~base-anclaje-dado de concreto. En particular, resulta imperativo cuantificar si el espesor
requerido de la placa $t_{req}$ es sensible a los cambios normativos, y en qué medida los factores
$R_{et}$ y $R_{ev}$ modifican la resistencia al cono de concreto respecto a la metodología
anterior. La presente tesina aborda esta brecha mediante un análisis paramétrico sistemático y
comparativo.

% ============================================================
% SECCIÓN 2.2 — OBJETIVO GENERAL
% ============================================================
\section{Objetivo General}

Realizar un análisis comparativo cuantitativo del diseño de placas base para columnas de acero
bajo la metodología de las NTC-DCEA 2017 y las NTC-DCEA 2023, evaluando el impacto técnico de
los cambios normativos en los parámetros de diseño del sistema placa base-anclaje-dado de concreto
bajo un régimen de flexocompresión con gran excentricidad, mediante la implementación de
algoritmos paramétricos en Python.

% ============================================================
% SECCIÓN 2.3 — OBJETIVOS ESPECÍFICOS
% ============================================================
\section{Objetivos Específicos}

\begin{enumerate}
  \item Implementar computacionalmente los algoritmos de diseño de placas base conforme a la
    NTC-DCEA 2017 y la NTC-DCEA 2023, aplicados al caso de flexocompresión con gran excentricidad
    ($e > e_{crit}$), verificando los estados límite de la placa base, los pernos de anclaje y
    el sistema de anclaje al cono de concreto, con el perfil IR~305$\times$74.5 (A992,
    $F_y = 3{,}515\ \mathrm{kg/cm^2}$) como perfil de referencia y una placa base de
    50$\times$50~cm.

  \item Evaluar mediante un barrido paramétrico el efecto del factor de confinamiento por estribos
    en tensión $R_{et}$ en función de la profundidad de empotramiento efectivo $h_{ef}$ y la
    densidad del refuerzo transversal del dado (diámetro y espaciado de estribos), determinando la
    eficiencia de aprovechamiento $\eta = T_{ua}/R_t$ para ambas normas.

  \item Cuantificar el impacto geométrico del cambio en el requisito de distancia mínima al borde
    libre del pedestal ---de $6\,d_o \geq 100\ \mathrm{mm}$ (NTC-DCEA 2017) a
    $7\,d_o \geq 100\ \mathrm{mm}$ para aceros A325/A449 (NTC-DCEA 2023)--- sobre las dimensiones
    mínimas del dado de concreto, en función del diámetro nominal del ancla $d_o$.

  \item Determinar si los cambios normativos de la NTC-DCEA 2023 producen variaciones en el
    espesor requerido de la placa base $t_{req}$ cuando el momento último $M_u$ varía en el rango
    de 0 a $14\ \mathrm{ton{\cdot}m}$, identificando el umbral de excentricidad a partir del cual
    se observan diferencias cuantificables entre ambas normas.
\end{enumerate}

La Figura~\ref{fig:casos_estudio} sintetiza los cuatro escenarios de solicitación que definen el
espacio de carga abarcado por los algoritmos de diseño implementados.

\begin{figure}[H]
  \centering
  \includegraphics[width=0.78\textwidth]{\detokenize{figuras/009_Casos de estudio..png}}
  \caption{Escenarios de solicitación estructural estudiados: (1)~compresión axial pura,
    (2)~flexocompresión de pequeña excentricidad, (3)~flexocompresión con gran excentricidad
    ---caso de activación de anclas en tensión--- y (4)~cortante dominante
    \cite{Mahmoud_2021}.}
  \label{fig:casos_estudio}
\end{figure}
